% !TEX root = ../../tir_monograph_v12.tex
\chapter{Minimal Finite Cell and Tetrahedral Closure}
\label{ch:v12_tetrahedral_closure}

\section{Minimal finite isotropy}

Let
\[
\mathbf n_1,\ldots,\mathbf n_m\in\mathbb R^3,
\qquad
|\mathbf n_a|=1,
\]
be equal-weight local relation directions in the carrier constructed in
Chapter~\ref{ch:v12_euclidean_dimension}.  Define
\[
\mathbf M:=\sum_{a=1}^{m}\mathbf n_a,
\qquad
Q:=\sum_{a=1}^{m}\mathbf n_a\mathbf n_a^T.
\]
The finite isotropy conditions are
\[
\boxed{\mathbf M=0},
\qquad
\boxed{Q=\frac m3 I_3}.
\]
The second condition has rank three.  For \(m\le2\) the relation span is
rank-deficient.  If \(m=3\), the zero-moment equation places the third vector in
the span of the first two, so the rank is at most two.  Hence
\[
\boxed{m\ge4.}
\]

For \(m=4\), place the four vectors into
\[
N=(\mathbf n_1\ \mathbf n_2\ \mathbf n_3\ \mathbf n_4).
\]
Then
\[
NN^T=\frac43I_3,
\qquad
N\mathbf1=0.
\]
The Gram matrix \(G=N^TN\) is therefore
\[
\boxed{
G
=
\frac43I_4-\frac13\mathbf1\mathbf1^T,
}
\]
so
\[
\boxed{
\mathbf n_a\cdot\mathbf n_b
=
-\frac13
\quad(a\neq b).
}
\]
The minimal equal-weight finite full-isotropy cell is consequently the regular
tetrahedral frame.

A convenient exact representative is
\[
\mathbf n_1=\frac1{\sqrt3}(1,1,1),\quad
\mathbf n_2=\frac1{\sqrt3}(1,-1,-1),
\]
\[
\mathbf n_3=\frac1{\sqrt3}(-1,1,-1),\quad
\mathbf n_4=\frac1{\sqrt3}(-1,-1,1).
\]

\section{Exact shape data}

For distinct vertices,
\[
\|\mathbf n_a-\mathbf n_b\|^2
=
2-2\left(-\frac13\right)
=
\frac83,
\]
hence the normalized edge length is
\[
\boxed{\hat a=\sqrt{\frac83}}.
\]
The Euclidean volume of the regular tetrahedron is
\[
\boxed{
\hat V_{\Delta^3}
=
\frac{\hat a^3}{6\sqrt2}
=
\frac{8}{9\sqrt3}.
}
\]

\section{Independent qubit informational-completeness route}

A generic qubit density operator has three independent real Bloch coordinates.
An \(m\)-outcome normalized probability vector has at most \(m-1\) independent
real entries.  Informational completeness therefore requires
\[
m-1\ge3,
\qquad
\boxed{m\ge4}.
\]

Using the same tetrahedral directions, define
\[
P_a=\frac12(I+\mathbf n_a\cdot\boldsymbol\sigma),
\qquad
E_a=\frac12P_a
=\frac14(I+\mathbf n_a\cdot\boldsymbol\sigma).
\]
Since \(\sum_a\mathbf n_a=0\),
\[
\boxed{\sum_{a=1}^{4}E_a=I}.
\]
For \(a\neq b\),
\[
\boxed{\operatorname{Tr}(P_aP_b)=\frac13}.
\]
For
\[
\rho=\frac12(I+\mathbf r\cdot\boldsymbol\sigma)
\]
the outcome probabilities are
\[
p_a=\operatorname{Tr}(\rho E_a)
=\frac14(1+\mathbf r\cdot\mathbf n_a).
\]
Using
\[
\sum_a\mathbf n_a\mathbf n_a^T=\frac43I_3
\]
gives the exact reconstruction
\[
\boxed{
\mathbf r
=
3\sum_{a=1}^{4}p_a\mathbf n_a.
}
\]
Thus the minimal symmetric informationally complete qubit frame and the minimal
finite isotropic relation frame carry the same regular tetrahedral Gram data.

\section{Congruence-class closure}

Let \(N\) be an ordered spatial tetrahedral frame and \(M\) an ordered SIC
tetrahedral frame, with
\[
N^TN=M^TM=G,
\qquad
NN^T=MM^T=\frac43I_3.
\]
Define
\[
\boxed{Q_{\Delta}:=\frac34MN^T.}
\]
Then
\[
Q_{\Delta}N=M,
\qquad
Q_{\Delta}Q_{\Delta}^T=I_3,
\]
so
\[
\boxed{Q_{\Delta}\in O(3).}
\]
After compatible orientation of the ordered frames the representative may be
chosen in \(SO(3)\).  The two routes therefore occupy one exact tetrahedral
orthogonal-congruence class in
\(\operatorname{Herm}_0(2)\cong\mathbb R^3\).

The common promoted object at this stage is this Gram/congruence class and its
orthogonal shape invariants.  Semantic role binding between an informational
outcome label and a spatial adjacency label remains a separate downstream
binding.

\section{Fubini--Study shape crosswalk}

The Bloch central angle between tetrahedral directions obeys
\[
\cos\chi=-\frac13.
\]
For one spherical tetrahedral face the spherical cosine law gives the interior
angle
\[
\alpha=\frac{2\pi}{3}.
\]
The spherical excess on the unit Bloch sphere is \(\pi\).  Since the qubit
Fubini--Study metric satisfies
\[
ds_{FS}^2=\frac14ds_{S^2}^2,
\]
one face has
\[
\boxed{a_{FS}^{face}=\frac\pi4},
\]
and all four faces carry
\[
\boxed{a_{FS}^{tet}=\pi}.
\]
Therefore the dimensionless dual-shape coefficient is
\[
\boxed{
C_{\Delta/FS}
=
\frac{\hat V_{\Delta^3}}{a_{FS}^{tet}}
=
\frac{8}{9\sqrt3\,\pi}.
}
\]

\section{Legacy tetrahedral semantics}

The historical Appendix~\ref{app:tetrahedron} carries the sector-specific
operator labels
\(\{\hat I,\hat M,\hat A_\mu,\hat A_\tau\}\) and its own six edge assignments.
In v12 that object remains sector-level provenance.  The canonical geometric
tetrahedron of this chapter is the regular Gram class derived above; any
role-binding map between the two tetrahedral descriptions must preserve their
typed edge data and receives its own source receipt.

\section{Source and validation receipt}

The canonical theorem surfaces are
\begin{itemize}
\item \path{TIR/foundations/TIR_MINIMAL_ISOTROPIC_TETRAHEDRAL_CELL_V0_1.md};
\item \path{TIR/foundations/TIR_QUBIT_TETRAHEDRAL_INFORMATIONAL_COMPLETENESS_V0_1.md};
\item \path{TIR/integration/TIR_TETRAHEDRAL_CONGRUENCE_CLASS_CLOSURE_V0_1.md}.
\end{itemize}
The deterministic v12 receipt is generated by
\[
\path{TIR/validation/tir_v12_ch06_tetrahedral_closure_v0_1.py}.
\]
